\chapter{练习类型}\label{s:exercises}

每个好木匠都有一套螺丝刀，
每个好老师都有不同种类的练习，
用来检查学习者到底在学什么、
帮他们练习新技能、
并让他们保持投入。
本章先描述几种可以用来检查教学是否有效的练习，
然后看自动评分的最新技术，
最后探讨讨论、项目，以及其他需要更多人工作判断来评估的重要工作类型。
我们的讨论部分取材于
\hreffoot{http://web-cat.org/questionbank/}{坎特伯雷题库}~\cite{Sand2013}，
它为入门计算的各种语言和主题提供了条目。

\seclbl{经典题型}{s:exercises-classics}

正如\secref{s:models-formative-assessment}里讨论的，
当错误答案能探查具体误解时，\emph{多项选择题}（MCQ）最有效。\index{多项选择题}
它们通常被设计来测试布鲁姆分类法的较低层次\index{布鲁姆分类法}
（\secref{s:process-objectives}），
但也可以要求学习者运用判断力。

\begin{aside}{一道多项选择题}
  当计算机求表达式 \texttt{price\ =\ addTaxes(cost\ -\ discount)} 的值时，
  各运算按什么顺序发生？
  \begin{enumerate}
  \item
    减法、函数调用、赋值
  \item
    函数调用、减法、赋值
  \item
    函数调用，然后同时进行赋值和减法
  \item
    以上都不对
  \end{enumerate}
\end{aside}

第二类经典的编程练习是\emph{写并运行}（C\&R），
学习者编写代码，产生指定的输出。
C\&R 练习可以像老师想要的那么简单或复杂，
但在课堂上使用时，
它们应该简短，并且只有一两个看似合理的正确答案。
对新手来说，让他们计算并打印一个值、
或调用一个特定函数，往往就够了：
有经验的老师常常忘了，
弄清哪个参数该放哪儿可能有多难。
对更进阶的学习者来说，
弄清该调用哪个函数更有趣，
也更能衡量他们的理解。

\begin{aside}{写并运行}
  变量 \texttt{picture} 装着从文件读入的一幅全彩图像。
  用同一个函数，
  创建这幅图像的黑白版本，
  并把它赋给一个叫 \texttt{monochrome} 的新变量。
\end{aside}

写并运行的练习可以和 MCQ 结合。
例如，
下面这道 MCQ 只有运行 Unix 的 \texttt{ls} 命令才能作答：

\begin{aside}{把 MCQ 与写并运行结合}
  你处在目录 \texttt{/home} 下。
  下列文件里，哪一个\emph{不}在该目录中？

  \begin{enumerate}
  \item
    \texttt{autumn.csv}
  \item
    \texttt{fall.csv}
  \item
    \texttt{spring.csv}
  \item
    \texttt{winter.csv}
  \end{enumerate}
\end{aside}

C\&R 能帮人练习他们最想学的技能，
但它们很难评估：
得到正确答案的方式可能有很多出乎意料，
而且如果自动评分系统因为代码和老师的不一样就拒绝它，
人们会很受打击。
减少这种情况发生的一种办法，是只评估他们的输出，
但这样就没法反馈他们在编程方式上的问题。
另一种办法是给他们一套小测试，让他们在提交前自己跑一跑
（提交时再对照一套更全面的测试运行）。
这有助于他们在做任何自认为会影响成绩的事之前，
弄清自己是不是完全误解了练习的意图。

除了编写满足某些规格的代码，
还可以让学习者写测试，
来判断一段代码是否符合规格。
这本身就是一项有用的技能，
而且这么做也能让学习者对老师有多辛苦多一点共情。

\begin{aside}{把写并运行反过来}
  函数 \texttt{monotonic\_sum} 计算一个数字列表中每一段的和，
  这些段里的值严格递增。
  例如，
  给定输入 \texttt{[1,\ 3,\ 3,\ 4,\ 5,\ 1]}，
  输出是 \texttt{[4,\ 12,\ 1]}。
  写并运行单元测试，判断下列哪个 bug 是该函数包含的：

  \begin{itemize}
  \item
    把每个负数都当作一个新子序列的开始。
  \item
    子和不包括每个子序列的第一个值。
  \item
    子和不包括每个子序列的最后一个值。
  \item
    只在值下降时（而不是在值未能递增时）重新开始求和。
  \end{itemize}
\end{aside}

\emph{填空}是 C\&R 的一种细化，
学习者拿到一些起始代码，必须把它补全。
（实践中，大多数 C\&R 练习其实都是填空，
因为老师会提供注释，
提醒学习者该采取哪些步骤。）
这类题目是渐隐示例的基础；
正如\chapref{s:architecture}里讨论的，
新手往往觉得它们比从头写全部代码更不吓人，
而且既然老师已经给出了答案的大部分结构，
提交的内容就可预测得多，因此也更容易检查。

\begin{aside}{填空}
  填空，
  让下面的代码打印字符串 \texttt{'hat'}。

\begin{minted}{text}
text = 'all that it is'
slice = text[____:____]
print(slice)
\end{minted}
\end{aside}

帕森斯问题\index{帕森斯问题}也避免了「令人恐惧的空白屏幕」问题，
同时让学习者能把控制流和词汇分开来集中练习
\cite{Pars2006,Eric2015,Morr2016,Eric2017}。
在线构建和完成帕森斯问题的工具有现成的~\cite{Ihan2011}，
但也可以通过让学习者在编辑器里重排代码行来模拟
（尽管有点笨拙）。

\begin{aside}{帕森斯问题}
  重排并缩进下面这些行，来求一个列表中正数之和。
  （你还需要在合适的地方加上冒号。）

\begin{minted}{text}
total = 0
if v > 0
total += v
for v in values
\end{minted}
\end{aside}

注意，给学习者的行数多于他们所需的行数，
或让他们重排一些行、再加几行，
都会让帕森斯问题显著变难~\cite{Harm2016}。

\seclbl{追踪}{s:exercises-tracing}

\emph{追踪执行}是帕森斯问题的逆过程：
给定几行代码，
学习者必须追踪这些行的执行顺序。
这既是一项基本的调试技能，\index{调试}
也是巩固学习者对循环、条件、
以及函数和方法调用求值顺序理解的好办法。
最简单的实现方式，是让学习者写出一串带标注的步骤。
让他们从一组选项里选出正确的序列
（也就是把它做成 MCQ）
会增加认知负荷却不增加价值，
因为他们得先费劲弄出正确序列，
再到选项列表里去找它。

\begin{aside}{追踪执行顺序}
  这个代码块里带标注的各行，按什么顺序执行？

\begin{minted}{text}
A)     vals = [-1, 0, 1]
B)     inverse_sum = 0
       try:
           for v in vals:
C)             inverse_sum += 1/v
       except:
D)         pass
\end{minted}
\end{aside}

\emph{追踪取值}和追踪执行类似，
但不是列出代码执行的顺序，
而是由学习者列出程序运行时一个或多个变量所取的值。
一种实现方式，是给学习者一张表，
它的列标注变量名、
行标注行号，
让他填出这些行上变量所取的值。

\newpage
\begin{aside}{追踪取值}
  这个程序执行时，\texttt{left} 和 \texttt{right} 各取什么值？

\begin{minted}{text}
A) left = 23
B) right = 6
C) while right:
D)     left, right = right, left % right
\end{minted}
\end{aside}

\begin{center}
\begin{tabular}{|l|l|l|}
  \hline
  行 & \texttt{left} & \texttt{right} \\
  \hline
  & & \\
  \hline
  & & \\
  \hline
  & & \\
  \hline
  & & \\
  \hline
  & & \\
  \hline
  & & \\
  \hline
\end{tabular}
\end{center}

你也可以要求学习者反向追踪代码，
弄清要产生某个特定结果，输入必须是什么~\cite{Armo2008}。
这类\emph{反向执行}问题需要搜索和演绎推理，
而当输出是一条错误信息时，
它们能帮学习者发展宝贵的调试技能。

\begin{aside}{反向执行}
  填出 \texttt{values} 里缺失的那个数，
  它导致了这个函数崩溃。

\begin{minted}{text}
values = [ [1.0, -0.5], [3.0, 1.5], [2.5, ___] ]
runningTotal = 0.0
for (reading, scaling) in values:
    runningTotal += reading / scaling
\end{minted}
\end{aside}

\emph{最小修改}练习也能帮学习者发展调试技能。
给定几行包含 bug 的代码，
学习者必须找到它，并做一处小改动来修复。
做改动可以用 C\&R，
而定位它可以用多项选择题。

\begin{aside}{最小修改}
  这个函数本该测试
  一个数是否落在某个范围内。
  做一处小改动，让它真正能做到。

\begin{minted}{text}
def inside(point, lower, higher):
    if (point <= lower):
        return false
    elif (point <= higher):
        return false
    else:
        return true
\end{minted}
\end{aside}

\emph{主题与变奏}练习类似，
但要求学习者做一处小改动，以某种特定方式改变输出，
而不是为修 bug 做改动。
允许的改动可以包括改变变量的初值、
把一次函数调用换成另一次、
交换内外层循环、
或改变复杂条件里各测试的顺序。
再说一遍，
这类练习给了学习者练习一项有用的真实技能的机会：
产出他们所需代码的最快方式，
就是微调一段已经能把事情做成的代码。

\begin{aside}{主题与变奏}
  改动下面函数里的内层循环，
  让它把图像的左上三角形填成指定的颜色。

\begin{minted}{text}
function fillTriangle(picture, color) is
    for x := 1 to picture.width do
        for y := 1 to picture.height do
            picture[x, y] = color
        end
    end
end
\end{minted}
\end{aside}

\emph{重构}练习是主题与变奏练习的补充：
给定一段能工作的代码，
学习者必须在不\emph{不}改变其输出的前提下以某种方式修改它。
例如，
学习者可以把循环换成向量化表达式，
或化简 while 循环里的条件。
这也是一项有用的真实技能，
但重构代码的方式往往太多，
所以评分需要人工检查。

\begin{aside}{重构}
  写一个列表推导式，
  效果和这个循环相同。

\begin{minted}{text}
result = []
for v in values:
    if len(v) > threshold:
        result.append(v)
\end{minted}
\end{aside}

\seclbl{图表}{s:exercises-diagrams}

让学习者画概念图和其他图表，
能让人洞察他们是怎么思考的（\secref{s:memory-concept-maps}），
但自由形式的图表需要人工花时间和判断来评估。
而\emph{标注图表}，
在教学上几乎同样有力，
却容易规模化得多。

与其让学习者从头创建图表，
不如给他们一张图和一组标签，
让他们把标签放到图上正确的位置。
图可以是数据结构
（「这段代码执行之后，哪些变量指向这个结构的哪些部分？」）、
一张图表（「把每一段代码和它生成的那部分图表匹配起来」），
或代码本身（「把每个术语和该程序元素的一个例子匹配起来」）。

\begin{aside}{标注图表}
  \figref{f:exercises-labeling}展示了
  一小段 HTML 在内存里是如何表示的。
  把标签 1—9 放到树的各个元素上，
  以显示在深度优先遍历中到达它们的顺序。
\end{aside}

\figpdf{figures/labeling.pdf}{标注图表}{f:exercises-labeling}{}

使用图表的另一种办法，
是给学习者图表的各个部件，让他们正确地摆出来。
这是帕森斯问题在视觉上的等价物，
你可以视他们准备好与否，提供或多或少的骨架来帮他们定位。
我至今怀念当年为了在某个点得到正确电压、
而试着在电路图里摆放电阻和电容的经历；
也见过老师给学习者一组固定的 Scratch 积木，\index{Scratch}
让他们只用这些积木画出某个特定图形。

\emph{匹配题}可以看作标注题的一个特例，
其中「图」是一列文字，
标签则取自另一列。
\emph{一对一匹配}给学习者两份等长的列表，
让他们把对应的项配对，
例如「把每段代码和它产生的输出匹配起来」。

\begin{aside}{匹配}
  把\figref{f:exercises-matching}里每个正则表达式运算符
  和它的作用匹配起来。
\end{aside}

\figpdfhere{figures/matching.pdf}{匹配题}{f:exercises-matching}{}

在\emph{多对多匹配}里，两份列表长度不等，
所以有些项可能匹配到好几个，
另一些则可能一个也匹配不上。
多对多更难，
因为学习者没法先做容易的匹配来缩小搜索空间。
匹配题可以通过让学习者提交配对列表来实现
（比如「A3, B1, C2」），
但那很笨拙、容易出错。
让他们在 MCQ 里识别一组正确配对更糟，
因为看错实在太容易了。
画或拖拽要好得多，
但可能需要花些功夫来实现。

\emph{排序}是匹配题的一个特例，
它（稍微）更适合通过列表来作答，
因为我们的头脑相当擅长发现序列里的错误或异常。
排序标准决定了所需的推理层次。
如果你让学习者把各种排序算法从最快排到最慢，
你大概考的是回忆
（也就是要求他们认出算法名字、知道它们的性质）；
而让他们把各种解从最稳健排到最脆弱，
考的则是推理和判断。

\emph{概括}也要求学习者运用高阶思维，
并给他们机会练习一项报告 bug 时非常有用的技能。
例如，你可以问学习者：
「哪一句话最能描述当 x 从 0 变到 10 时 f 的输出如何变化？」
然后用一道多项选择题给出几个选项。
你也可以在限定领域里要求非常简短的自由作答，
比如「稳定排序算法的关键特征是什么？」
不用人工，我们没法为这些问题做全自动检查而不产生多得令人沮丧的
假阳性（接受错误答案）
和假阴性（拒绝正确答案），
但这类问题很适合同伴评分（\secref{s:individual-peer}）。

\seclbl{自动评分}{s:exercises-grading}
\index{自动评分}

自动程序评分工具存在的时间比我活的年头还长：
最早的公开记载可以追溯到 1960 年~\cite{Holl1960}，
而~\cite{Douc2005,Ihan2010}里发表的综述按名字提到了许多具体工具。
构建这类工具远比乍看起来复杂。
作业是怎么表示的？
提交是如何跟踪和报告的？
学习者能协作吗？
提交如何安全地执行？
\cite{Edwa2014a}整整一篇论文都在讲
一种自适应检测和管理代码提交里死循环的方案，
而这还只是冒出来的众多问题之一。

讨论自动评分器时，
重要的是把学习者满意度和学习成效区分开。
例如，
\cite{Magu2018}把一门二年级 CS 课的非正式编程实验
换成使用自动评分器的每周机评测验。
学习者不喜欢这套自动化系统，
但整门课的不及格率减半了，
拿到一等荣誉的学习者人数增加了两倍。
相反，
\cite{Rubi2014}也开始使用一套为竞赛设计的自动评分器，
却没看到学习者的退课率有显著下降；
再一次，
学习者对工具发表了一些负面评论，
作者把这归因于它反馈信息的质量，
而不是因为他们不喜欢自动评分。

\cite{Srid2016}采取了另一种做法。
他们用\gref{g:fuzz-testing}{fuzz testing}
（也就是随机生成的测试用例）
来检查学习者代码是否和老师提供的参考实现做同样的事。
在一门 1400 人的入门课程的第一个项目里，
fuzz testing 抓住了手写测试用例套件漏掉的错误，
涉及超过 48\% 的学习者。

\cite{Basu2015}给学习者一套解答测试用例，
但学习者必须先回答关于其预期行为的问题，才能解锁每一个用例、
从而把它用在自己的解上。
例如，
假设学习者要写一个函数，找出列表里相邻数字之和最大的一对。
在被允许使用该题的测试之前，
他们必须为「\texttt{largestPair(4, 3, -1, 5, 3, 3)} 会产出什么？」
选出正确答案。
在一门 1300 人的大学课程里，
绝大多数学习者都选择用这种方式先确认自己对测试用例的理解，
再动手解题，
随后他们提问更少，对作业的困惑也更少。

\begin{aside}{反对现成工具}
  用现成的风格检查工具来给学习者代码打分很诱人。
  然而，
  \cite{Nutb2016}起初发现，人工给的分数
  和风格检查器报出的规则违反之间没有相关性。
  有时这是因为学习者多次违反同一条规则
  （因此丢的分比该丢的多），
  但另一些时候是因为他们提交的是几乎没有改动的作业起始代码，
  却拿到了比该得的分更多的分数。

  即便是专门为教学构建的工具，也可能达不到老师的需要。
  \cite{Keun2016a,Keun2016b}考察了 69 个自动评分工具产生的信息。
  他们发现，这些工具常常不给出关于如何修问题、如何走下一步的反馈。
  他们还发现，大多数老师没法轻易地把这些工具改成自己需要的：
  像许多工作流工具一样，
  它们往往把创造者对机构如何运作的、连自己都没意识到的假设强加于人。
  他们的分类方案，是挑选这类工具时一份有用的购物清单。
\end{aside}

\cite{Buff2015}对提供自动反馈的整个想法做了见识广博的反思。
他们的出发点是：
「自动评分系统能帮学习者找出代码里的 bug，
[但]可能无意中让学习者不愿批判性地思考和充分地测试，
反而鼓励他们依赖老师的测试。」
他们识别出的关键问题之一是，
学习者可能已经充分地测试了自己的代码，
但某个功能可能仍然没有按老师的规格实现。
在这种情况下，
「失败」不是因为没有测试，
而是因为误解了要求，
而且更多的测试也不太可能暴露这个问题。
如果自动评分系统不提供有洞见、可行动的反馈，
这段经历只会让学习者受挫。

为了提供那种反馈，
\cite{Buff2015}的系统会找出学习者代码里哪些方法被失败的测试执行过，
这样系统就能把失败的测试
和学习者提交里特定的功能关联起来。
系统通过看学习者是否已经充分测试了相关功能，
来判断具体提示是否已经「挣得」，
所以学习者没法靠提示来替代做测试。

\cite{Srid2016}描述了其他一些在自动测试学习者代码时与他们分享反馈的办法。
第一种是提供测试的预期输出——但
那样学习者就会为这些输入写死输出
（因为任何能被钻空子的东西都会被钻）。
第二种是报告学习者代码的通过/失败结果，
但只在提交截止日期之后才提供测试的实际输入和输出。
然而，
告诉学习者他们错了却不告诉他们为什么，会让人沮丧。

第三种做法是用一种叫\gref{g:hashing}{哈希}的技巧，
生成一个依赖输出、却不泄露输出的值。
如果用户产出的输出完全正确，
它的哈希就能解锁答案，
但没法从哈希反推出输出应该是什么。
哈希需要更多工作来解释和搭建，
但它在过早给出答案
和在该给出答案时却不给出之间，取得了不错的平衡。

\seclbl{高阶思维}{s:exercises-higher}

许多其他种类的编程练习，
在超过一小撮学习者的班里，老师就很难评估，
对自动平台来说更难评估。
更大的编程项目（但愿）是课程正在指向的目标，
但给出反馈的唯一方式就是逐个来。

\emph{代码评阅}一般来说也难以自动评分，
但如果给学习者一份要查找的缺陷清单，
并让他们把具体的评语和具体的代码行匹配起来，就可以处理。
例如，可以告诉学习者有两处缩进错误和一处糟糕的变量名，
让他们指出来。
如果他们更进阶，
可以给他们半打关于这段代码能做出的评论种类，
而不告诉他们每种各该找到几处。

\cite{Steg2016b}是代码风格评分标准的一个好起点，
而~\cite{Luxt2009}更一般地考察了编程课里的同伴评阅。
如果你要让学习者做评阅，
就用校准式同伴互评（\secref{s:individual-peer}），\index{校准式同伴互评}
好让他们有「好的反馈该长什么样」的范本。

\begin{aside}{代码评阅}
  用提供的评分标准标出每行代码里的问题。

\begin{minted}{text}
01)  def addem(f):
02)      x1 = open(f).readlines()
03)      x2 = [x for x in x1 if x.strip()]
04)      changes = 0
05)      for v in x2:
06)          print('total', total)
07)          tot = tot + int(v)
08)      print('total')
\end{minted}

  \begin{longtable}{ll}
    1. 变量名差   & 2. 使用了未定义的变量 \\
    3. 缺少返回值 & 4. 未使用的变量
  \end{longtable}
\end{aside}

\seclbl{练习}{s:exercises-exercises}

\exercise{写并运行}{两人一组}{10}

创建一道简短的 C\&R 练习，
然后和同伴交换，
看看你们各自要多久才能看懂并做完对方的练习。
练习描述里有歧义或误解吗？

\exercise{把写并运行反过来}{小组}{15}

分成 4—6 人的小组。
让组里每个成员创建一道反向 C\&R 练习，
要求人们弄清什么输入能产生某个特定输出。
随机挑两道，
看看全组能找到多少个满足要求的、各不相同的输入。

\exercise{追踪取值}{两人一组}{10}

写一个简短的程序（10—15 行），
和同伴交换，
并追踪程序里各个变量随时间如何变化取值。
你和同伴记录追踪过程的方式有什么不同？

\exercise{重构}{小组}{15}

分成 3—4 人的小组。
让每个人选一段自己写过的、还不够整洁的短代码（10—30 行），
然后随机挑一段，让组里每个人各自把它整理干净。
你们整理后的版本有什么不同？
如果自动评分、或在大班里评分，
你应付得了所有这些变体的程度有多好或多差？

\exercise{标注图表}{两人一组}{10}

画一张能展示你最近讲过的东西的图：
浏览器如何从服务器取数据、
对象和类之间的关系、
或者 R 里数据框是如何被索引的。
把标签放在旁边，让同伴来放。

\exercise{纸笔谜题}{全班}{15}

\cite{Butl2017}描述了一组纸笔谜题，
它们可以变成入门编程作业，
并报告说这些作业被学习者喜爱、也能鼓励元认知。
想一个你小时候玩过的简单纸笔谜题或游戏，
描述你会怎样把它变成一道编程练习。

\exercise{数一数失败}{两人一组}{15}

任何对练习需要多少时间的有用估计，
都必须把失败出现的频率、
以及失败损失的时间考虑进去。
例如，
编辑文本文件看起来是个简单任务，
但找到这些文件呢？
大多数 GUI 编辑器把东西存到用户的桌面或主目录；
如果课程用到的文件存在别处，
相当一部分人没有帮助就没法导航到正确的目录。
（如果你觉得这是个小问题，
请重新读一读\chapref{s:memory}里关于专家盲区的讨论。）

和同伴一起，
列出你用过或做过的练习里见过的那些出岔子的「简单」事。
它们多久出现一次？
学习者自己或靠帮助修好它们要多久？
你现在在课堂上为应对它们留了多少时间预算？

\exercise{说到时间}{个人}{10}

到目前为止，本书里练习的时间估计有多准？
